% Part: lambda-calculus
% Chapter: lambda-definability
% Section: fixpoints

\documentclass[../../../include/open-logic-section]{subfiles}

\begin{document}

\olfileid{lam}{ldf}{fp}

\olsection{स्थिरबिंदू}

क्रमगुणित फलनाचे पुनरावर्तनाने असे पद $\fn{Fac}$
म्हणून परिभाषण करायचे आहे असे समजा:
\[
\fn{Fac} \ident \lambd[n][\fn{IsZero}\, n\, \num 1 (\fn{Mult}\, n (\fn{Fac}(\fn{Pred} \, n)))]
\]
म्हणजे $n = 0$ असेल, तर $n$ चे क्रमगुणित $1$;
अन्यथा ते $n$ गुणिले $n-1$ चे क्रमगुणित.
पण $\fn{Fac}$ ची अशी व्याख्या करता येत नाही,
कारण उजव्या बाजूस $\fn{Fac}$ हेच पद येते.
अशा स्वसंदर्भाचा समावेश असलेल्या पुनरावर्ती
व्याख्या लॅम्डा कलनाचा भाग नाहीत. उदाहरणार्थ,
पुढील पद परिभाषित करताना
\[
\fn{Mult} \ident \lambd[ab][b (\fn{Add}\, a) \num{0}]
\]
उजव्या बाजूस $\fn{Add}$ यांसारखी पूर्वी परिभाषित
पदेच वापरली जातात. $\fn{Add}$ च्या जागी त्याचे
परिभाषक पद ठेवून ते काढून टाकता येते. मग
$\fn{Mult}$ शुद्ध लॅम्डा पद म्हणून मिळेल;
$\fn{Add}$ मध्येही इतर परिभाषित पदे असतील
(उदा., $\fn{Add}'$ मध्ये आहेत), तर अशी जागापालट
सुरू ठेवून शेवटी शुद्ध लॅम्डा पद मिळवता येते.

वरच्या $\fn{Fac}$ सारख्या पुनरावर्ती व्याख्येच्या
बाबतीत मात्र तसे होत नाही. उजवीकडील $\fn{Fac}$
च्या जागी $\fn{Fac}$ चीच व्याख्या ठेवल्यास पुढील पद मिळते:
\begin{align*}
  B &\ident \lambd[m][\fn{IsZero}\,m\,\num{1}
       (\fn{Mult}\,m\,(\fn{Fac}(\fn{Pred}\,m)))]\\
  \fn{Fac} &\ident \lambd[n][\fn{IsZero}\,n\,\num{1}
       (\fn{Mult}\,n\,(B(\fn{Pred}\,n)))]
\end{align*}
येथे B हे फक्त आतले पद थोडक्यात लिहिण्यासाठीचे नाव आहे.
त्याच्या व्याख्येतील $\fn{Fac}$ अजूनही गेलेले नाही.
बाह्य अमूर्तीकरणाची व्याप्ती दोन्ही conditional शाखांवर आहे.
ही प्रक्रिया पुन्हा केल्यास व्याख्या वाढतच राहते
आणि कधीही शुद्ध लॅम्डा पद मिळत नाही. म्हणून
क्रमगुणित, किंवा सर्वसाधारणपणे पुनरावर्ती फलने,
अशा रीतीने परिभाषित करता येत नाहीत.

तरी या पुनरावर्ती व्याख्येतून एक गोष्ट समजते:
$f$ हे क्रमगुणित फलन दर्शवणारे पद असेल, तर
\[
\fn{Fac}' \ident \lambd[g][\lambd[n][\fn{IsZero} \, n \, \num 1\, (\fn{Mult} \, n\, (g (\fn{Pred} n)))]]
\]
या पदाला $f$ लागू केल्यावर मिळणारे
$\fn{Fac}'\,f$ हेही क्रमगुणित फलन दर्शवते.
म्हणजे $\fn{Fac}'$ याकडे फलन स्वीकारून फलन
परत करणारे फलन म्हणून पाहिल्यास, $\fn{Fac'}\, f$
आणि~$f$ यांचे सर्व चर्च अंकांवरील परिणाम समान असतात.
यावरून ही दोन पदे स्वतः बीटा-सममूल्य आहेत असे मात्र
ठरत नाही; संख्यात्मक फलनाचे निरूपण आणि पदांची
बीटा-सममूल्यता या वेगळ्या अटी आहेत. $\fn{Fac}' \, f \equal[\beta] f$
हा गुणधर्म असलेल्या फलन~$f$ ला $\fn{Fac}'$
चा \emph{स्थिरबिंदू} म्हणतात. पुढील संयोजकाने
$\fn{Fac}'$ चा असा स्थिरबिंदू मिळवून त्याचे सर्व
अंकांवरील परिणाम क्रमगुणिताचे आहेत हे सिद्ध करू.

लॅम्डा कलनात दिलेल्या पदाचा स्थिरबिंदू काढणारी
पदे आहेत. त्यांच्या साहाय्याने $\fn{Fac}'$
सारख्या पदापासून क्रमगुणिताची व्याख्या मिळवता येते.

\begin{defn}
  \ollabel{defn:Turing-Y}
  \emph{Y-संयोजक} हे पुढील पद आहे:
  \[
  Y \ident (\lambd[ux][x(uux)])(\lambd[ux][x(uux)]).
  \]
\end{defn}

\begin{thm}
  $Y$ साठी कोणत्याही पद $g$ बाबत $Yg \red g(Yg)$; म्हणून
  $Yg$ हा नेहमी~$g$ चा स्थिरबिंदू असतो.
\end{thm}

\begin{proof}
  $(\lambd[ux][x(uux)])$ ला~$U$ असे संक्षेपाने
  लिहू. मग $Y \ident UU$. म्हणून
  \begin{align*}
    Y g &\ident (\lambd[ux][x(uux)])U\,g \\
    &\red (\lambd[x][x(UUx)])g\\
    & \red g(UUg) \ident g(Yg).
  \end{align*}
  $g(Yg)$ आणि $Yg$ या दोघांचेही $g(Yg)$ मध्ये
  न्यूनीकरण होते; म्हणून $g(Yg) \equal[\beta]
  Yg$ आणि $Yg$ हा~$g$ चा स्थिरबिंदू आहे.
\end{proof}

$Yg$ मध्ये रेडेक्स असल्याने न्यूनीकरण अनंतकाळ
पुढे चालू राहू शकते:
\begin{align*}
  Y g &\red g (Y g) \\
    &\red g (g (Y g)) \\
    &\red g(g (g (Y g)))\\
    &\ldots
\end{align*}
हा सांत पदांच्या अधिकाधिक उकलांचा अंतहीन अनुक्रम आहे;
तो स्वतः एक अनंत लॅम्डा पद नाही. म्हणून $Yg$ म्हणजे $g$
वारंवार लावलेल्या उकलांनी व्यक्त केलेला स्थिरबिंदू असे
सहज अर्थाने समजू शकतो. त्याला $g$ आणखी एकदा
लावल्याने, बोलायचे तर, काहीच जास्त होत नाही:
$Yg$ वर $g$ अनंत वेळा लावल्यावर पुन्हा $g$
लावले, तरी ते $Yg$ वर $g$ अनंत वेळा लावण्यासारखेच आहे.

वरील $Yg$ पासून सुरू होणारा $\beta$-न्यूनीकरणाचा
अनुक्रम अनंत आहे, हे लक्षात घ्या. त्यामुळे $Yg$
एखाद्या पदाला लागू केल्यावरचे $(Yg)N$ हे पदही
$(g(Yg))N$, $(g(g(Yg)))N$, \dots अशा अनंत
भिन्न पदांमध्ये न्यूनीत होऊ शकते. तरी न्यूनीकरणाचा
एखादा \emph{दुसरा} अनुक्रम सामान्य रूपात संपणे
शक्य आहे.

क्रमगुणिताचे उदाहरण घ्या. $\fn{Fac}$ ची व्याख्या
$Y\, \fn{Fac}'$ अशी करा (म्हणजे $\fn{Fac'}$
चा स्थिरबिंदू). मग:
\begin{align*}
  \fn{Fac}\,\num 3
  &\red Y\, \fn{Fac}' \, \num 3 \\
  &\red \fn{Fac}' (Y \,\fn{Fac}') \, \num 3 \\
  &\ident (\lambd[x][\lambd[n][\fn{IsZero} \, n\, \num 1\,
      (\fn{Mult}\, n\, (x (\fn{Pred}\, n)))]]) \, \fn{Fac} \, \num 3\\
  & \red \fn{IsZero} \, \num 3\, \num 1\,
      (\fn{Mult}\, \num 3\, (\fn{Fac} (\fn{Pred}\, \num 3))) \\
  &\red  \fn{Mult}\, \num 3 \, (\fn{Fac} \, \num 2).
  \intertext{त्याचप्रमाणे,}
  \fn{Fac}\,\num 2
  &\red  \fn{Mult}\, \num 2 \, (\fn{Fac} \, \num 1) \\
  \fn{Fac}\,\num 1
  &\red  \fn{Mult}\, \num 1 \, (\fn{Fac} \, \num 0)
  \intertext{पण}
  \fn{Fac}\,\num 0
  &\red \fn{Fac}' (Y \,\fn{Fac}') \, \num 0 \\
  &\ident (\lambd[x][\lambd[n][\fn{IsZero} \, n\, \num 1\,
      (\fn{Mult}\, n\, (x (\fn{Pred}\, n)))]]) \, \fn{Fac} \, \num 0\\
  & \red \fn{IsZero} \, \num 0\, \num 1\,
      (\fn{Mult}\, \num 0\, (\fn{Fac} (\fn{Pred}\, \num 0))).\\
  &\red \num 1.
  \intertext{म्हणून एकत्रितपणे}
  \fn{Fac}\,\num 3
  &\red  \fn{Mult}\, \num 3 \,
  (\fn{Mult} \, \num 2\, (\fn{Mult} \, \num 1\, \num 1)).
\end{align*}

$\fn{Fac'}$ साठी जे खरे आहे ते कोणत्याही
पुनरावर्ती व्याख्येसाठी खरे आहे. पुढील पुनरावर्ती
समीकरण दिले आहे असे समजा. पुनरावर्तनाचे नाव आणि
सर्व निविष्ट प्राचलांची नावे परस्पर भिन्न निवडायची आहेत:
\begin{align*}
g\,x_1\dots x_n & \equal[\beta] N
\intertext{येथे $N$ मध्ये~$g$ आणि $x_1$, \dots,~$x_n$
असू शकतात. मग असे पद~$G \ident (Y \lambd[g][\lambd[x_1\dots x_n][N]])$
नेहमी असते की}
G\,x_1 \dots x_n & \equal[\beta] \Subst{N}{G}{g}.
\intertext{कारण स्थिरबिंदू प्रमेयानुसार,}
G \ident (Y \lambd[g][\lambd[x_1\dots x_n][N]]) & \red \lambd[g][\lambd[x_1\dots x_n][N]](Y \lambd[g][\lambd[x_1\dots x_n][N]])\\
& \ident (\lambd[g][\lambd[x_1\dots x_n][N]])\,G
\intertext{आणि परिणामी,}
G\,x_1 \dots x_n
& \red (\lambd[g][\lambd[x_1\dots x_n][N]])\,G\,x_1\dots x_n\\
& \red (\lambd[x_1\dots x_n][\Subst{N}{G}{g}])\,x_1\dots x_n\\
  & \red \Subst{N}{G}{g}.
\end{align*}

\olref{defn:Turing-Y} मधील $Y$-संयोजक
ॲलन ट्यूरिंगने दिला होता. ॲलॉन्झो चर्चने
दुसरे रूप सुचवले होते; त्याला~$Y_C$ म्हणू:
\[
Y_C \ident \lambd[g][(\lambd[x][g(xx)])(\lambd[x][g(xx)])].
\]
चर्चचा संयोजक ट्यूरिंगच्या संयोजकापेक्षा काहीसा
कमकुवत आहे. औपचारिक चर g वापरल्यास $Y_Cg \equal[\beta] g(Y_Cg)$, पण
$Y_Cg \bred g(Y_Cg)$ नाही. $V$ हे
$\lambd[x][g(xx)]$ पद असू द्या; येथे त्याचे बद्ध नाव
आदेशित पदाच्या मुक्त चरांपासून वेगळे निवडायचे आहे. मग
$Y_C \ident \lambd[g][VV]$. त्यामुळे
\begin{align*}
  VV & \ident (\lambd[x][g(xx)])V \red g(VV)
  \text{ आणि त्यामुळे}\\
  Y_C g & \ident (\lambd[g][VV])g \red VV \red g(VV), \text{ पण तसेच}\\
  g(Y_C g) & \ident g((\lambd[g][VV])g) \red g(VV).
\end{align*}
दुसऱ्या शब्दांत, $Y_Cg$ आणि $g(Y_Cg)$ यांचे
न्यूनीकरण होऊन एकच पद $g(VV)$ मिळते;
म्हणून $Y_Cg \equal[\beta] g(Y_Cg)$. अनेक
उपयोजनांसाठी एवढे पुरेसे असते.

\end{document}
